\documentclass[../../main.tex]{subfiles}
\begin{document}
\section{习题}\label{section:ch05-习题}

\begin{exercise}\label{exercise:ch05-p5-1}
设 $M$ 是(不带边的)光滑流形, $S\subseteq M$ 是浸入子流形。证明: $S$ 在 $M$ 中闭,当且仅当包含映射 $S\hookrightarrow M$ 是固有映射。
\end{exercise}

\begin{exercise}\label{exercise:ch05-p5-2}
证明\cref{theorem:ch05-boundary-submanifold}:若 $M$ 是带边光滑流形,则 $\partial M$ 是 $M$ 的真嵌入子流形。
\end{exercise}

\begin{exercise}\label{exercise:ch05-p5-3}
设 $M\subseteq\mathbb{R}^n$ 是嵌入子流形。证明 $M$ 是闭子集,当且仅当它在 $M$ 上的包含映射是固有映射。
\end{exercise}

\begin{exercise}\label{exercise:ch05-p5-4}
设 $S\subseteq\mathbb{R}^3$ 是旋转曲面,由曲线 $C\subseteq H=\{(r,z):r>0\}$ 生成。证明 $C$ 是 $S$ 到 $(r,z)$ 平面上投影的逆像。
\end{exercise}

\begin{exercise}\label{exercise:ch05-p5-5}
设 $F:M\to N$ 是光滑淹没, $S\subseteq N$ 是浸入子流形。证明 $F^{-1}(S)$ 是 $M$ 的浸入子流形,其余维数等于 $S$ 的余维数。
\end{exercise}

\begin{exercise}\label{exercise:ch05-p5-6}
设 $M$ 是(不带边的)光滑流形, $S\subseteq M$ 是浸入子流形, $S'$ 是 $S$ 的任一开子集。证明 $S'$ 是 $M$ 的浸入子流形,并说明它一般不是嵌入子流形。
\end{exercise}

\begin{exercise}\label{exercise:ch05-p5-7}
设 $S\subseteq M$ 是浸入子流形。证明 $S$ 在 $M$ 中闭,当且仅当包含映射 $S\hookrightarrow M$ 是固有映射。
\end{exercise}

\begin{exercise}\label{exercise:ch05-p5-8}
设 $M$ 是光滑流形, $S\subseteq M$ 是嵌入子流形, $\Phi:M\to N$ 是光滑映射。证明:若 $\Phi|_S:S\to N$ 是光滑嵌入,则 $S$ 在 $M$ 中有邻域 $U$,使 $\Phi|_U:U\to N$ 是光滑嵌入。
\end{exercise}

\begin{exercise}\label{exercise:ch05-p5-9}
设 $S\subseteq M$ 是浸入子流形, $\iota:S\to M$ 为包含映射。证明:在 $S$ 上使 $\iota$ 为浸入的拓扑与光滑结构存在,当且仅当 $S$ 容许一个光滑流形结构使 $\iota$ 为光滑浸入。
\end{exercise}

\begin{exercise}\label{exercise:ch05-p5-10}
证明\cref{proposition:ch05-tangent-boundary}。
\end{exercise}

\begin{exercise}\label{exercise:ch05-p5-11}
证明\cref{figure:ch05-bad-level-sets} 中 $\Phi$ 与 $\Psi$ 的零点集都不是 $\mathbb{R}^2$ 的嵌入子流形。
\end{exercise}

\begin{exercise}\label{exercise:ch05-p5-12}
设 $S\subseteq M$ 是浸入子流形, $\iota:S\to M$ 为包含映射。若 $S$ 紧,证明 $\iota$ 是光滑嵌入。
\end{exercise}

\begin{exercise}\label{exercise:ch05-p5-13}
证明\cref{proposition:ch05-level-sets-boundary}。
\end{exercise}

\begin{exercise}\label{exercise:ch05-p5-14}
完成\cref{theorem:ch05-slice-boundary} 的证明。
\end{exercise}

\begin{exercise}\label{exercise:ch05-p5-15}
证明\cref{example:ch05-hemisphere-submanifold} 中的 $D$ 是 $S^2$ 的带边嵌入子流形。
\end{exercise}

\end{document}